\documentclass[12pt]{article}
\usepackage[T1]{fontenc}
\usepackage[utf8]{inputenc}
\usepackage{lmodern}
\usepackage{textcomp}
\usepackage{enumitem}
\usepackage{geometry}
\usepackage{graphicx}
\usepackage{amsmath, amssymb}
\geometry{margin=1in}
\begin{document}
\begin{center}
Economics - Graduate Level Assessment 10 \
Total Marks: 40 \
Answer all questions.
\end{center}

\section*{Multiple Choice (10 marks)}
\begin{enumerate}[label=\arabic*.]
\item How does an increasing national debt impact the market for U.S. dollars and the value of the dollar with respect to other currencies? MARKET FOR THE DOLLAR VALUE OF THE DOLLAR
\begin{enumerate}[label=(\alph*)]
    \item Increased supply Depreciating
    \item Decreased supply Appreciating
    \item No change in supply Appreciating
    \item Increased demand Depreciating
    \item Decreased demand Appreciating
\end{enumerate}
\item Which of the following would NOT shift the aggregate demand curve?
\begin{enumerate}[label=(\alph*)]
    \item A change in spending by state governments.
    \item A change in the inflation rate.
    \item A change in technology.
    \item A change in consumer confidence.
    \item A change in foreign trade policies.
\end{enumerate}
\item If Real GDP = \$200 billion and the price index = 200 Nominal GDP is
\begin{enumerate}[label=(\alph*)]
    \item \$300 billion
    \item \$800 billion
    \item \$4 billion
    \item \$600 billion
    \item \$200 billion
\end{enumerate}
\item The concentration ratio for a monopoly is
\begin{enumerate}[label=(\alph*)]
    \item 50
    \item 5
    \item 10
    \item 90
    \item 15
\end{enumerate}
\item What is meant by the term "normal goods"?
\begin{enumerate}[label=(\alph*)]
    \item Normal goods are those whose demand is independent of income.
    \item Normal goods are those whose demand curve shifts downward with an increase in income.
    \item Normal goods are those whose consumption is deemed socially unacceptable as income rises.
    \item Normal goods are those whose demand curve shifts upward with a decrease in income.
    \item Normal goods are those whose demand curve varies directly with income.
\end{enumerate}
\end{enumerate}

\section*{True / False (10 marks)}
\textit{Answer True or False}
\begin{enumerate}[label=\arabic*.]
\setcounter{enumi}{5}
\item True or False: The residual from a standard regression model is the difference between the actual value, y, and the fitted value, y-hat.
\item True or False: In the short run, a perfectly competitive firm may operate at a loss depending on market prices relative to average total cost.
\item True or False: A negative current account balance necessitates a decrease in the country's interest rates to stabilize the economy.
\item True or False: The measure of a firm's monopoly pricing power is represented by the formula (P-MC)/$MC^2$, indicating greater power with larger values.
\item True or False: A university degree is classified as capital because it directly generates income like physical assets do.
\end{enumerate}

\section*{Long Form Response (20 marks)}
\textit{Answer each question with a clear and complete explanation.}
\begin{enumerate}[label=\arabic*.]
\setcounter{enumi}{10}
\item Discuss the implications of falling prices and rising real GDP in an economy, focusing on how technological advancements could drive these trends and their potential effects on economic stability.
\item Discuss the implications of a reserve ratio of 0.1 on a bank's lending capacity, specifically analyzing how a $200 deposit can lead to a maximum lending potential of $2,000.
\end{enumerate}

\vfill
\noindent\textit{End of Paper}
\end{document}